\documentclass[10pt]{article}
\usepackage{tmlr}
%%%%% NEW MATH DEFINITIONS %%%%%

\usepackage{amsmath,amsfonts,bm}

% Mark sections of captions for referring to divisions of figures
\newcommand{\figleft}{{\em (Left)}}
\newcommand{\figcenter}{{\em (Center)}}
\newcommand{\figright}{{\em (Right)}}
\newcommand{\figtop}{{\em (Top)}}
\newcommand{\figbottom}{{\em (Bottom)}}
\newcommand{\captiona}{{\em (a)}}
\newcommand{\captionb}{{\em (b)}}
\newcommand{\captionc}{{\em (c)}}
\newcommand{\captiond}{{\em (d)}}

% Highlight a newly defined term
\newcommand{\newterm}[1]{{\bf #1}}


% Figure reference, lower-case.
\def\figref#1{figure~\ref{#1}}
% Figure reference, capital. For start of sentence
\def\Figref#1{Figure~\ref{#1}}
\def\twofigref#1#2{figures \ref{#1} and \ref{#2}}
\def\quadfigref#1#2#3#4{figures \ref{#1}, \ref{#2}, \ref{#3} and \ref{#4}}
% Section reference, lower-case.
\def\secref#1{section~\ref{#1}}
% Section reference, capital.
\def\Secref#1{Section~\ref{#1}}
% Reference to two sections.
\def\twosecrefs#1#2{sections \ref{#1} and \ref{#2}}
% Reference to three sections.
\def\secrefs#1#2#3{sections \ref{#1}, \ref{#2} and \ref{#3}}
% Reference to an equation, lower-case.
\def\eqref#1{equation~\ref{#1}}
% Reference to an equation, upper case
\def\Eqref#1{Equation~\ref{#1}}
% A raw reference to an equation---avoid using if possible
\def\plaineqref#1{\ref{#1}}
% Reference to a chapter, lower-case.
\def\chapref#1{chapter~\ref{#1}}
% Reference to an equation, upper case.
\def\Chapref#1{Chapter~\ref{#1}}
% Reference to a range of chapters
\def\rangechapref#1#2{chapters\ref{#1}--\ref{#2}}
% Reference to an algorithm, lower-case.
\def\algref#1{algorithm~\ref{#1}}
% Reference to an algorithm, upper case.
\def\Algref#1{Algorithm~\ref{#1}}
\def\twoalgref#1#2{algorithms \ref{#1} and \ref{#2}}
\def\Twoalgref#1#2{Algorithms \ref{#1} and \ref{#2}}
% Reference to a part, lower case
\def\partref#1{part~\ref{#1}}
% Reference to a part, upper case
\def\Partref#1{Part~\ref{#1}}
\def\twopartref#1#2{parts \ref{#1} and \ref{#2}}

\def\ceil#1{\lceil #1 \rceil}
\def\floor#1{\lfloor #1 \rfloor}
\def\1{\bm{1}}
\newcommand{\train}{\mathcal{D}}
\newcommand{\valid}{\mathcal{D_{\mathrm{valid}}}}
\newcommand{\test}{\mathcal{D_{\mathrm{test}}}}

\def\eps{{\epsilon}}


% Random variables
\def\reta{{\textnormal{$\eta$}}}
\def\ra{{\textnormal{a}}}
\def\rb{{\textnormal{b}}}
\def\rc{{\textnormal{c}}}
\def\rd{{\textnormal{d}}}
\def\re{{\textnormal{e}}}
\def\rf{{\textnormal{f}}}
\def\rg{{\textnormal{g}}}
\def\rh{{\textnormal{h}}}
\def\ri{{\textnormal{i}}}
\def\rj{{\textnormal{j}}}
\def\rk{{\textnormal{k}}}
\def\rl{{\textnormal{l}}}
% rm is already a command, just don't name any random variables m
\def\rn{{\textnormal{n}}}
\def\ro{{\textnormal{o}}}
\def\rp{{\textnormal{p}}}
\def\rq{{\textnormal{q}}}
\def\rr{{\textnormal{r}}}
\def\rs{{\textnormal{s}}}
\def\rt{{\textnormal{t}}}
\def\ru{{\textnormal{u}}}
\def\rv{{\textnormal{v}}}
\def\rw{{\textnormal{w}}}
\def\rx{{\textnormal{x}}}
\def\ry{{\textnormal{y}}}
\def\rz{{\textnormal{z}}}

% Random vectors
\def\rvepsilon{{\mathbf{\epsilon}}}
\def\rvtheta{{\mathbf{\theta}}}
\def\rva{{\mathbf{a}}}
\def\rvb{{\mathbf{b}}}
\def\rvc{{\mathbf{c}}}
\def\rvd{{\mathbf{d}}}
\def\rve{{\mathbf{e}}}
\def\rvf{{\mathbf{f}}}
\def\rvg{{\mathbf{g}}}
\def\rvh{{\mathbf{h}}}
\def\rvu{{\mathbf{i}}}
\def\rvj{{\mathbf{j}}}
\def\rvk{{\mathbf{k}}}
\def\rvl{{\mathbf{l}}}
\def\rvm{{\mathbf{m}}}
\def\rvn{{\mathbf{n}}}
\def\rvo{{\mathbf{o}}}
\def\rvp{{\mathbf{p}}}
\def\rvq{{\mathbf{q}}}
\def\rvr{{\mathbf{r}}}
\def\rvs{{\mathbf{s}}}
\def\rvt{{\mathbf{t}}}
\def\rvu{{\mathbf{u}}}
\def\rvv{{\mathbf{v}}}
\def\rvw{{\mathbf{w}}}
\def\rvx{{\mathbf{x}}}
\def\rvy{{\mathbf{y}}}
\def\rvz{{\mathbf{z}}}

% Elements of random vectors
\def\erva{{\textnormal{a}}}
\def\ervb{{\textnormal{b}}}
\def\ervc{{\textnormal{c}}}
\def\ervd{{\textnormal{d}}}
\def\erve{{\textnormal{e}}}
\def\ervf{{\textnormal{f}}}
\def\ervg{{\textnormal{g}}}
\def\ervh{{\textnormal{h}}}
\def\ervi{{\textnormal{i}}}
\def\ervj{{\textnormal{j}}}
\def\ervk{{\textnormal{k}}}
\def\ervl{{\textnormal{l}}}
\def\ervm{{\textnormal{m}}}
\def\ervn{{\textnormal{n}}}
\def\ervo{{\textnormal{o}}}
\def\ervp{{\textnormal{p}}}
\def\ervq{{\textnormal{q}}}
\def\ervr{{\textnormal{r}}}
\def\ervs{{\textnormal{s}}}
\def\ervt{{\textnormal{t}}}
\def\ervu{{\textnormal{u}}}
\def\ervv{{\textnormal{v}}}
\def\ervw{{\textnormal{w}}}
\def\ervx{{\textnormal{x}}}
\def\ervy{{\textnormal{y}}}
\def\ervz{{\textnormal{z}}}

% Random matrices
\def\rmA{{\mathbf{A}}}
\def\rmB{{\mathbf{B}}}
\def\rmC{{\mathbf{C}}}
\def\rmD{{\mathbf{D}}}
\def\rmE{{\mathbf{E}}}
\def\rmF{{\mathbf{F}}}
\def\rmG{{\mathbf{G}}}
\def\rmH{{\mathbf{H}}}
\def\rmI{{\mathbf{I}}}
\def\rmJ{{\mathbf{J}}}
\def\rmK{{\mathbf{K}}}
\def\rmL{{\mathbf{L}}}
\def\rmM{{\mathbf{M}}}
\def\rmN{{\mathbf{N}}}
\def\rmO{{\mathbf{O}}}
\def\rmP{{\mathbf{P}}}
\def\rmQ{{\mathbf{Q}}}
\def\rmR{{\mathbf{R}}}
\def\rmS{{\mathbf{S}}}
\def\rmT{{\mathbf{T}}}
\def\rmU{{\mathbf{U}}}
\def\rmV{{\mathbf{V}}}
\def\rmW{{\mathbf{W}}}
\def\rmX{{\mathbf{X}}}
\def\rmY{{\mathbf{Y}}}
\def\rmZ{{\mathbf{Z}}}

% Elements of random matrices
\def\ermA{{\textnormal{A}}}
\def\ermB{{\textnormal{B}}}
\def\ermC{{\textnormal{C}}}
\def\ermD{{\textnormal{D}}}
\def\ermE{{\textnormal{E}}}
\def\ermF{{\textnormal{F}}}
\def\ermG{{\textnormal{G}}}
\def\ermH{{\textnormal{H}}}
\def\ermI{{\textnormal{I}}}
\def\ermJ{{\textnormal{J}}}
\def\ermK{{\textnormal{K}}}
\def\ermL{{\textnormal{L}}}
\def\ermM{{\textnormal{M}}}
\def\ermN{{\textnormal{N}}}
\def\ermO{{\textnormal{O}}}
\def\ermP{{\textnormal{P}}}
\def\ermQ{{\textnormal{Q}}}
\def\ermR{{\textnormal{R}}}
\def\ermS{{\textnormal{S}}}
\def\ermT{{\textnormal{T}}}
\def\ermU{{\textnormal{U}}}
\def\ermV{{\textnormal{V}}}
\def\ermW{{\textnormal{W}}}
\def\ermX{{\textnormal{X}}}
\def\ermY{{\textnormal{Y}}}
\def\ermZ{{\textnormal{Z}}}

% Vectors
\def\vzero{{\bm{0}}}
\def\vone{{\bm{1}}}
\def\vmu{{\bm{\mu}}}
\def\vtheta{{\bm{\theta}}}
\def\va{{\bm{a}}}
\def\vb{{\bm{b}}}
\def\vc{{\bm{c}}}
\def\vd{{\bm{d}}}
\def\ve{{\bm{e}}}
\def\vf{{\bm{f}}}
\def\vg{{\bm{g}}}
\def\vh{{\bm{h}}}
\def\vi{{\bm{i}}}
\def\vj{{\bm{j}}}
\def\vk{{\bm{k}}}
\def\vl{{\bm{l}}}
\def\vm{{\bm{m}}}
\def\vn{{\bm{n}}}
\def\vo{{\bm{o}}}
\def\vp{{\bm{p}}}
\def\vq{{\bm{q}}}
\def\vr{{\bm{r}}}
\def\vs{{\bm{s}}}
\def\vt{{\bm{t}}}
\def\vu{{\bm{u}}}
\def\vv{{\bm{v}}}
\def\vw{{\bm{w}}}
\def\vx{{\bm{x}}}
\def\vy{{\bm{y}}}
\def\vz{{\bm{z}}}

% Elements of vectors
\def\evalpha{{\alpha}}
\def\evbeta{{\beta}}
\def\evepsilon{{\epsilon}}
\def\evlambda{{\lambda}}
\def\evomega{{\omega}}
\def\evmu{{\mu}}
\def\evpsi{{\psi}}
\def\evsigma{{\sigma}}
\def\evtheta{{\theta}}
\def\eva{{a}}
\def\evb{{b}}
\def\evc{{c}}
\def\evd{{d}}
\def\eve{{e}}
\def\evf{{f}}
\def\evg{{g}}
\def\evh{{h}}
\def\evi{{i}}
\def\evj{{j}}
\def\evk{{k}}
\def\evl{{l}}
\def\evm{{m}}
\def\evn{{n}}
\def\evo{{o}}
\def\evp{{p}}
\def\evq{{q}}
\def\evr{{r}}
\def\evs{{s}}
\def\evt{{t}}
\def\evu{{u}}
\def\evv{{v}}
\def\evw{{w}}
\def\evx{{x}}
\def\evy{{y}}
\def\evz{{z}}

% Matrix
\def\mA{{\bm{A}}}
\def\mB{{\bm{B}}}
\def\mC{{\bm{C}}}
\def\mD{{\bm{D}}}
\def\mE{{\bm{E}}}
\def\mF{{\bm{F}}}
\def\mG{{\bm{G}}}
\def\mH{{\bm{H}}}
\def\mI{{\bm{I}}}
\def\mJ{{\bm{J}}}
\def\mK{{\bm{K}}}
\def\mL{{\bm{L}}}
\def\mM{{\bm{M}}}
\def\mN{{\bm{N}}}
\def\mO{{\bm{O}}}
\def\mP{{\bm{P}}}
\def\mQ{{\bm{Q}}}
\def\mR{{\bm{R}}}
\def\mS{{\bm{S}}}
\def\mT{{\bm{T}}}
\def\mU{{\bm{U}}}
\def\mV{{\bm{V}}}
\def\mW{{\bm{W}}}
\def\mX{{\bm{X}}}
\def\mY{{\bm{Y}}}
\def\mZ{{\bm{Z}}}
\def\mBeta{{\bm{\beta}}}
\def\mPhi{{\bm{\Phi}}}
\def\mLambda{{\bm{\Lambda}}}
\def\mSigma{{\bm{\Sigma}}}

% Tensor
\DeclareMathAlphabet{\mathsfit}{\encodingdefault}{\sfdefault}{m}{sl}
\SetMathAlphabet{\mathsfit}{bold}{\encodingdefault}{\sfdefault}{bx}{n}
\newcommand{\tens}[1]{\bm{\mathsfit{#1}}}
\def\tA{{\tens{A}}}
\def\tB{{\tens{B}}}
\def\tC{{\tens{C}}}
\def\tD{{\tens{D}}}
\def\tE{{\tens{E}}}
\def\tF{{\tens{F}}}
\def\tG{{\tens{G}}}
\def\tH{{\tens{H}}}
\def\tI{{\tens{I}}}
\def\tJ{{\tens{J}}}
\def\tK{{\tens{K}}}
\def\tL{{\tens{L}}}
\def\tM{{\tens{M}}}
\def\tN{{\tens{N}}}
\def\tO{{\tens{O}}}
\def\tP{{\tens{P}}}
\def\tQ{{\tens{Q}}}
\def\tR{{\tens{R}}}
\def\tS{{\tens{S}}}
\def\tT{{\tens{T}}}
\def\tU{{\tens{U}}}
\def\tV{{\tens{V}}}
\def\tW{{\tens{W}}}
\def\tX{{\tens{X}}}
\def\tY{{\tens{Y}}}
\def\tZ{{\tens{Z}}}


% Graph
\def\gA{{\mathcal{A}}}
\def\gB{{\mathcal{B}}}
\def\gC{{\mathcal{C}}}
\def\gD{{\mathcal{D}}}
\def\gE{{\mathcal{E}}}
\def\gF{{\mathcal{F}}}
\def\gG{{\mathcal{G}}}
\def\gH{{\mathcal{H}}}
\def\gI{{\mathcal{I}}}
\def\gJ{{\mathcal{J}}}
\def\gK{{\mathcal{K}}}
\def\gL{{\mathcal{L}}}
\def\gM{{\mathcal{M}}}
\def\gN{{\mathcal{N}}}
\def\gO{{\mathcal{O}}}
\def\gP{{\mathcal{P}}}
\def\gQ{{\mathcal{Q}}}
\def\gR{{\mathcal{R}}}
\def\gS{{\mathcal{S}}}
\def\gT{{\mathcal{T}}}
\def\gU{{\mathcal{U}}}
\def\gV{{\mathcal{V}}}
\def\gW{{\mathcal{W}}}
\def\gX{{\mathcal{X}}}
\def\gY{{\mathcal{Y}}}
\def\gZ{{\mathcal{Z}}}

% Sets
\def\sA{{\mathbb{A}}}
\def\sB{{\mathbb{B}}}
\def\sC{{\mathbb{C}}}
\def\sD{{\mathbb{D}}}
% Don't use a set called E, because this would be the same as our symbol
% for expectation.
\def\sF{{\mathbb{F}}}
\def\sG{{\mathbb{G}}}
\def\sH{{\mathbb{H}}}
\def\sI{{\mathbb{I}}}
\def\sJ{{\mathbb{J}}}
\def\sK{{\mathbb{K}}}
\def\sL{{\mathbb{L}}}
\def\sM{{\mathbb{M}}}
\def\sN{{\mathbb{N}}}
\def\sO{{\mathbb{O}}}
\def\sP{{\mathbb{P}}}
\def\sQ{{\mathbb{Q}}}
\def\sR{{\mathbb{R}}}
\def\sS{{\mathbb{S}}}
\def\sT{{\mathbb{T}}}
\def\sU{{\mathbb{U}}}
\def\sV{{\mathbb{V}}}
\def\sW{{\mathbb{W}}}
\def\sX{{\mathbb{X}}}
\def\sY{{\mathbb{Y}}}
\def\sZ{{\mathbb{Z}}}

% Entries of a matrix
\def\emLambda{{\Lambda}}
\def\emA{{A}}
\def\emB{{B}}
\def\emC{{C}}
\def\emD{{D}}
\def\emE{{E}}
\def\emF{{F}}
\def\emG{{G}}
\def\emH{{H}}
\def\emI{{I}}
\def\emJ{{J}}
\def\emK{{K}}
\def\emL{{L}}
\def\emM{{M}}
\def\emN{{N}}
\def\emO{{O}}
\def\emP{{P}}
\def\emQ{{Q}}
\def\emR{{R}}
\def\emS{{S}}
\def\emT{{T}}
\def\emU{{U}}
\def\emV{{V}}
\def\emW{{W}}
\def\emX{{X}}
\def\emY{{Y}}
\def\emZ{{Z}}
\def\emSigma{{\Sigma}}

% entries of a tensor
% Same font as tensor, without \bm wrapper
\newcommand{\etens}[1]{\mathsfit{#1}}
\def\etLambda{{\etens{\Lambda}}}
\def\etA{{\etens{A}}}
\def\etB{{\etens{B}}}
\def\etC{{\etens{C}}}
\def\etD{{\etens{D}}}
\def\etE{{\etens{E}}}
\def\etF{{\etens{F}}}
\def\etG{{\etens{G}}}
\def\etH{{\etens{H}}}
\def\etI{{\etens{I}}}
\def\etJ{{\etens{J}}}
\def\etK{{\etens{K}}}
\def\etL{{\etens{L}}}
\def\etM{{\etens{M}}}
\def\etN{{\etens{N}}}
\def\etO{{\etens{O}}}
\def\etP{{\etens{P}}}
\def\etQ{{\etens{Q}}}
\def\etR{{\etens{R}}}
\def\etS{{\etens{S}}}
\def\etT{{\etens{T}}}
\def\etU{{\etens{U}}}
\def\etV{{\etens{V}}}
\def\etW{{\etens{W}}}
\def\etX{{\etens{X}}}
\def\etY{{\etens{Y}}}
\def\etZ{{\etens{Z}}}

% The true underlying data generating distribution
\newcommand{\pdata}{p_{\rm{data}}}
% The empirical distribution defined by the training set
\newcommand{\ptrain}{\hat{p}_{\rm{data}}}
\newcommand{\Ptrain}{\hat{P}_{\rm{data}}}
% The model distribution
\newcommand{\pmodel}{p_{\rm{model}}}
\newcommand{\Pmodel}{P_{\rm{model}}}
\newcommand{\ptildemodel}{\tilde{p}_{\rm{model}}}
% Stochastic autoencoder distributions
\newcommand{\pencode}{p_{\rm{encoder}}}
\newcommand{\pdecode}{p_{\rm{decoder}}}
\newcommand{\precons}{p_{\rm{reconstruct}}}

\newcommand{\laplace}{\mathrm{Laplace}} % Laplace distribution

\newcommand{\E}{\mathbb{E}}
\newcommand{\Ls}{\mathcal{L}}
\newcommand{\R}{\mathbb{R}}
\newcommand{\emp}{\tilde{p}}
\newcommand{\lr}{\alpha}
\newcommand{\reg}{\lambda}
\newcommand{\rect}{\mathrm{rectifier}}
\newcommand{\softmax}{\mathrm{softmax}}
\newcommand{\sigmoid}{\sigma}
\newcommand{\softplus}{\zeta}
\newcommand{\KL}{D_{\mathrm{KL}}}
\newcommand{\Var}{\mathrm{Var}}
\newcommand{\standarderror}{\mathrm{SE}}
\newcommand{\Cov}{\mathrm{Cov}}
% Wolfram Mathworld says $L^2$ is for function spaces and $\ell^2$ is for vectors
% But then they seem to use $L^2$ for vectors throughout the site, and so does
% wikipedia.
\newcommand{\normlzero}{L^0}
\newcommand{\normlone}{L^1}
\newcommand{\normltwo}{L^2}
\newcommand{\normlp}{L^p}
\newcommand{\normmax}{L^\infty}

\newcommand{\parents}{Pa} % See usage in notation.tex. Chosen to match Daphne's book.

\DeclareMathOperator*{\argmax}{arg\,max}
\DeclareMathOperator*{\argmin}{arg\,min}

\DeclareMathOperator{\sign}{sign}
\DeclareMathOperator{\Tr}{Tr}
\let\ab\allowbreak

\usepackage{hyperref}
\usepackage{url}
\usepackage{calc}
\usepackage{tikz}
\usetikzlibrary{positioning,arrows.meta,calc}
\usepackage{longtable,booktabs,array}
\providecommand{\tightlist}{\relax}
\usepackage{xeCJK}
\setCJKmainfont{Songti SC}

\title{分解树节点该得多少分：一个闭式、可负控制的条件信息度量（玩具模型示范）}
% 匿名投稿（TMLR 双盲）：不出现作者块。

\def\month{十一月}
\def\year{2026}

\begin{document}

\maketitle

\begin{abstract}
把一个判断拆成分解树之后，一个自然但没有统一答案的问题是：每个节点各贡献了多少信息？我们给出一个可闭式计算的条件信息量：把\emph{浅层判断} $S$（在等速率假设下按占比枚举）与\emph{深层判断} $D$（施加真实速率后得到）分开，用 $H(D|S)$ 度量深层相对浅层的\textbf{增量信息}。该量只依赖硬件规格与一个工作负载先验：不需要实验、不需要训练、不需要标注。在四个平台 × 40,000 个采样工作负载上：浅层判断有 43--56\% 的情形判错瓶颈类型，深层修正带来的时间占比移动为 0.17--0.23，深层独有信息 $H(D|S)$ 为 0.71--0.77 bits。两个负控制把两种平凡解释排除——"深层只是转发浅层信息"或"深层因定义而占优"：平坦速率对照下该增量\textbf{精确为零}，且在 24 种参数—深度置换下不变。三个手算锚点精确复现。在所有判据先于数据生成提交（预注册 \citep{camerer2018}）的前提下，本文把这些结果作为一个\textbf{评分基础}的示范，而非任何方法或对齐问题的解法。
\end{abstract}

\section{1. 为什么是这个量}\label{why-this-quantity}

\textbf{占比归因不够。} 瓶颈判断是由交互项决定的组合函数。在 Jetson Thor T5000 上，FP16 只占 70/30 负载的 30\%，施加速率后却占 63\% 的时间（50/50 时为 80\%）。占比小不等于贡献小；继承线性假设的度量（占比、固定上下文下的消融）会系统性低估交互项。

\textbf{用信息族度量而非占比。} 互信息无需线性假设即可捕捉非线性依赖：即便 $D$ 是 $S$ 的非线性函数，仍有 $I(S;D)>0$。分解 $H(D)=H(D|S)+I(S;D)$ 把"浅层已知后仍然存活的信息"单独隔离出来——这正是我们想给深层节点记的账。

\textbf{两层结构让"增量"可定义。} 浅层 = 等速率下的占比枚举；深层 = 施加真实速率 $t_i=f_i/\mathrm{rate}_i$ 后的判断。$H(D|S)$ 就是"给定浅层后深层新增多少"。

\textbf{操作化形式。} 深层以两种方式修正浅层：方向修正（瓶颈类型翻转，$S\neq D$）和粒度修正（时间占比移动）。翻转率与占比移动是可解释的分数量；$H(D|S)$ 是其信息论严格形式。

我们明确声明：链式法则与互信息是教科书信息论，roofline 是教科书计算机体系结构 \citep{roofline2009}，PID 的独有/冗余/协同分解是既有工作。我们不主张这些工具本身新。贡献在于\textbf{把它们组合成一个可闭式计算、可负控制、可预注册的评分协议 \citep{camerer2018}}，并在一个玩具域上完整检验它。


% Figure 1 (中文版): 评分协议三合一结构图
% 用法：在 paper-zh-tmlr.tex 的 §1 末尾 % Figure 1 (中文版): 评分协议三合一结构图
% 用法：在 paper-zh-tmlr.tex 的 §1 末尾 % Figure 1 (中文版): 评分协议三合一结构图
% 用法：在 paper-zh-tmlr.tex 的 §1 末尾 \input{fig1-scoring-protocol-zh}
% 依赖：\usepackage{tikz}, \usetikzlibrary{positioning,arrows.meta,calc}
% 说明：TikZ 逻辑与英文版 fig1-scoring-protocol.tex 完全一致，仅标签改为中文
\begin{figure}[t]
\centering
\begin{tikzpicture}[
  node distance=0.9cm and 1.1cm,
  box/.style={draw, rounded corners, align=center, minimum height=0.9cm, font=\small},
  lbl/.style={font=\scriptsize},
  arr/.style={->, >=Latex, thick, gray!70}
]
% 左列：分解树
\node[box, draw=blue!60, fill=blue!8] (root) {分解树\\(toy: 2-level)};
\node[box, below=of root] (s) {浅层判断 $S$\\(等速率·占比枚举)};
\node[box, below=of s] (d) {深层判断 $D$\\(真实速率 $t_i=f_i/r_i$)};

% 中列：量
\node[box, right=1.9cm of s, draw=red!60, fill=red!8] (inc) {增量信息\\$H(D|S)=H(D)-I(S;D)$};
\node[box, below=of inc] (score) {可解释分数\\flip $=P(S\neq D)$ · shift};

% 右列：对照
\node[box, right=2.0cm of inc, draw=green!60, fill=green!8] (flat) {平坦速率对照\\all $r_i=1$};
\node[box, below=of flat, draw=green!40, fill=green!5] (zero) {结果：\\$H(D|S)=0$ 精确};
\node[box, right=1.9cm of score, draw=orange!60, fill=orange!8] (perm) {24 参数—深度置换\\(多重集固定)};
\node[box, below=of perm, draw=orange!40, fill=orange!5] (inv) {结果：增量不变\\(去深度化)};

\draw[arr] (root) -- (s);
\draw[arr] (s) -- (d);
\draw[arr] (s.east) -- node[lbl, above] {条件化} (inc.west);
\draw[arr] (d.east) -- (inc.west);
\draw[arr] (inc) -- (score);
\draw[arr] (inc.east) -- node[lbl, above] {施于同一量} (flat.west);
\draw[arr] (flat) -- (zero);
\draw[arr] (score.east) -- (perm.west);
\draw[arr] (perm) -- (inv);
\end{tikzpicture}
\caption{评分协议：分解树 $\rightarrow$ 条件信息增量 $\rightarrow$ 两个负控制。左列给出 toy 结构；中列是闭式可算的评分量；右列两个臂分别排除``深层只是转发浅层''（平坦对照精确归零）与``深层因定义占优''（24 置换不变）。全部判据先于数据生成提交（预注册）。}
\label{fig:scoring-protocol}
\end{figure}

% 依赖：\usepackage{tikz}, \usetikzlibrary{positioning,arrows.meta,calc}
% 说明：TikZ 逻辑与英文版 fig1-scoring-protocol.tex 完全一致，仅标签改为中文
\begin{figure}[t]
\centering
\begin{tikzpicture}[
  node distance=0.9cm and 1.1cm,
  box/.style={draw, rounded corners, align=center, minimum height=0.9cm, font=\small},
  lbl/.style={font=\scriptsize},
  arr/.style={->, >=Latex, thick, gray!70}
]
% 左列：分解树
\node[box, draw=blue!60, fill=blue!8] (root) {分解树\\(toy: 2-level)};
\node[box, below=of root] (s) {浅层判断 $S$\\(等速率·占比枚举)};
\node[box, below=of s] (d) {深层判断 $D$\\(真实速率 $t_i=f_i/r_i$)};

% 中列：量
\node[box, right=1.9cm of s, draw=red!60, fill=red!8] (inc) {增量信息\\$H(D|S)=H(D)-I(S;D)$};
\node[box, below=of inc] (score) {可解释分数\\flip $=P(S\neq D)$ · shift};

% 右列：对照
\node[box, right=2.0cm of inc, draw=green!60, fill=green!8] (flat) {平坦速率对照\\all $r_i=1$};
\node[box, below=of flat, draw=green!40, fill=green!5] (zero) {结果：\\$H(D|S)=0$ 精确};
\node[box, right=1.9cm of score, draw=orange!60, fill=orange!8] (perm) {24 参数—深度置换\\(多重集固定)};
\node[box, below=of perm, draw=orange!40, fill=orange!5] (inv) {结果：增量不变\\(去深度化)};

\draw[arr] (root) -- (s);
\draw[arr] (s) -- (d);
\draw[arr] (s.east) -- node[lbl, above] {条件化} (inc.west);
\draw[arr] (d.east) -- (inc.west);
\draw[arr] (inc) -- (score);
\draw[arr] (inc.east) -- node[lbl, above] {施于同一量} (flat.west);
\draw[arr] (flat) -- (zero);
\draw[arr] (score.east) -- (perm.west);
\draw[arr] (perm) -- (inv);
\end{tikzpicture}
\caption{评分协议：分解树 $\rightarrow$ 条件信息增量 $\rightarrow$ 两个负控制。左列给出 toy 结构；中列是闭式可算的评分量；右列两个臂分别排除``深层只是转发浅层''（平坦对照精确归零）与``深层因定义占优''（24 置换不变）。全部判据先于数据生成提交（预注册）。}
\label{fig:scoring-protocol}
\end{figure}

% 依赖：\usepackage{tikz}, \usetikzlibrary{positioning,arrows.meta,calc}
% 说明：TikZ 逻辑与英文版 fig1-scoring-protocol.tex 完全一致，仅标签改为中文
\begin{figure}[t]
\centering
\begin{tikzpicture}[
  node distance=0.9cm and 1.1cm,
  box/.style={draw, rounded corners, align=center, minimum height=0.9cm, font=\small},
  lbl/.style={font=\scriptsize},
  arr/.style={->, >=Latex, thick, gray!70}
]
% 左列：分解树
\node[box, draw=blue!60, fill=blue!8] (root) {分解树\\(toy: 2-level)};
\node[box, below=of root] (s) {浅层判断 $S$\\(等速率·占比枚举)};
\node[box, below=of s] (d) {深层判断 $D$\\(真实速率 $t_i=f_i/r_i$)};

% 中列：量
\node[box, right=1.9cm of s, draw=red!60, fill=red!8] (inc) {增量信息\\$H(D|S)=H(D)-I(S;D)$};
\node[box, below=of inc] (score) {可解释分数\\flip $=P(S\neq D)$ · shift};

% 右列：对照
\node[box, right=2.0cm of inc, draw=green!60, fill=green!8] (flat) {平坦速率对照\\all $r_i=1$};
\node[box, below=of flat, draw=green!40, fill=green!5] (zero) {结果：\\$H(D|S)=0$ 精确};
\node[box, right=1.9cm of score, draw=orange!60, fill=orange!8] (perm) {24 参数—深度置换\\(多重集固定)};
\node[box, below=of perm, draw=orange!40, fill=orange!5] (inv) {结果：增量不变\\(去深度化)};

\draw[arr] (root) -- (s);
\draw[arr] (s) -- (d);
\draw[arr] (s.east) -- node[lbl, above] {条件化} (inc.west);
\draw[arr] (d.east) -- (inc.west);
\draw[arr] (inc) -- (score);
\draw[arr] (inc.east) -- node[lbl, above] {施于同一量} (flat.west);
\draw[arr] (flat) -- (zero);
\draw[arr] (score.east) -- (perm.west);
\draw[arr] (perm) -- (inv);
\end{tikzpicture}
\caption{评分协议：分解树 $\rightarrow$ 条件信息增量 $\rightarrow$ 两个负控制。左列给出 toy 结构；中列是闭式可算的评分量；右列两个臂分别排除``深层只是转发浅层''（平坦对照精确归零）与``深层因定义占优''（24 置换不变）。全部判据先于数据生成提交（预注册）。}
\label{fig:scoring-protocol}
\end{figure}


\section{2. 模型}\label{model}

工作负载由四个分量构成，占比 $f_i$ 从 $\mathrm{Dirichlet}(1,1,1,1)$ 采样。每个分量有速率 $\mathrm{rate}_i$。

\begin{itemize}
\tightlist
\item
  浅层判断 $S$：假设所有 $\mathrm{rate}_i=1$，则时间占比 $\propto f_i$，瓶颈 $=\arg\max_i f_i$。
\item
  深层判断 $D$：施加真实速率，$t_i=f_i/\mathrm{rate}_i$ \citep{rooflineTime2020}，瓶颈 $=\arg\max_i t_i$。
\end{itemize}

增量量：

\[
\text{flip}=P(S\neq D),\qquad \text{shift}=\mathbb{E}\,|\mathrm{share}_D-\mathrm{share}_S|,
\] \[
H(D|S)=H(D)-I(S;D).
\]

四个平台：RTX 5090（最大/最小速率比 51.9）、Thor T5000（9.9）、RTX 4090（287.4）、A100（44.3）。

\section{3. 结果}\label{results}

\begin{longtable}[]{@{}llll@{}}
\toprule\noalign{}
平台 & flip & shift & $H(D|S)$ (bits) \\
\midrule\noalign{}
\endhead
\bottomrule\noalign{}
\endlastfoot
RTX 5090 & 0.527 & 0.200 & 0.76 \\
Thor T5000 & 0.434 & 0.166 & 0.71 \\
RTX 4090 & 0.563 & 0.231 & 0.77 \\
A100 & 0.556 & 0.218 & 0.74 \\
\end{longtable}

三个手算锚点精确复现：T5000 50/50 → 80\%；T5000 70/30 → 63\%；5090 标量占 40\% 且发射速率砍半 → 慢 40\%。

修正量不随速率比的\emph{极端度}单调：A100（比 44.3）的 shift 高于 5090（比 51.9）。因此我们只主张与速率比\emph{结构}相关；在四个平台上这是 4 点结构观察，不是 scaling law。

\section{4. 负控制与防御}\label{negative-controls-and-defenses}

\textbf{（1）平坦速率对照——测得的臂。} 把占位硬件的所有速率设为 1。结果：flip = 0.0000，shift = 0.0000，$H(D|S)=0.0$——\textbf{精确}为零，而非仅仅很小。这排除"深层只是转发浅层信息"。

\textbf{（2）参数—深度置换——测得的臂。} 把速率向量在四个分量间做全部 24 种置换，固定速率多重集。若增量追踪的是"哪个分量被叫深层"这个\emph{标签}，它会随置换改变；若追踪的是速率比\emph{内容}，应不变。结果：恒等置换精确复现主结果；24 种置换下 flip 只在窄 Monte-Carlo 带内波动（5090：0.5251--0.5304）。判据被"去深度化"。

\textbf{（3）层内类型对照——设计层论证。} 在本文的两层模型里，浅层是枚举型、深层是判断型，两者只在速率向量上不同：类型不与深度混淆。早期的节点级模型确实混淆了二者，在那里类型对照\emph{反对}假设（枚举型 0.5841 bits > 判断型 0.1344 bits）。该数字来自另一个混淆的模型，不作为证据使用。

\textbf{HARKing 声明。} 全部判据、公式与生死线检验（真实平台上深层独有信息 > 0，平坦对照上 = 0）均在数据生成前提交（commit \texttt{eb1b0598c}，17:39，早于 \texttt{718394ece}，18:06）。唯一的事后改动\emph{弱化}了一条次级假设（随极端度单调 → 依赖结构），因为数据否证了强形式。

\section{5. 树结构示范（4 点观察）}\label{tree-demonstration-a-4-point-observation}

在两路切分（compute vs memory）下，修正量可分解为 branch-level 与 within-branch 两部分，四格之和为 1（加性自洽）：

\begin{longtable}[]{@{}
  >{\raggedright\arraybackslash}p{(\linewidth - 10\tabcolsep) * \real{0.1667}}
  >{\raggedright\arraybackslash}p{(\linewidth - 10\tabcolsep) * \real{0.1667}}
  >{\raggedright\arraybackslash}p{(\linewidth - 10\tabcolsep) * \real{0.1667}}
  >{\raggedright\arraybackslash}p{(\linewidth - 10\tabcolsep) * \real{0.1667}}
  >{\raggedright\arraybackslash}p{(\linewidth - 10\tabcolsep) * \real{0.1667}}
  >{\raggedright\arraybackslash}p{(\linewidth - 10\tabcolsep) * \real{0.1667}}@{}}
\toprule\noalign{}
\begin{minipage}[b]{\linewidth}\raggedright
平台
\end{minipage} & \begin{minipage}[b]{\linewidth}\raggedright
co-flip
\end{minipage} & \begin{minipage}[b]{\linewidth}\raggedright
within-flip
\end{minipage} & \begin{minipage}[b]{\linewidth}\raggedright
co-correct
\end{minipage} & \begin{minipage}[b]{\linewidth}\raggedright
within-correct
\end{minipage} & \begin{minipage}[b]{\linewidth}\raggedright
主导维度
\end{minipage} \\
\midrule\noalign{}
\endhead
\bottomrule\noalign{}
\endlastfoot
T5000 (9.9) & 0.530 & 0.145 & 0.036 & 0.290 & within \\
A100 (44.3) & 0.399 & 0.442 & 0.045 & 0.114 & branch \\
5090 (51.9) & 0.409 & 0.335 & 0.064 & 0.191 & branch \\
4090 (287.4) & 0.394 & 0.451 & 0.044 & 0.112 & branch \\
\end{longtable}

段级 flip 精确复现一维主结果。最温和的平台（T5000）是 within-branch 主导；三个极端平台是 branch 主导且内部聚簇。这是 4 点结构观察，不是连续权衡曲线；切分固定为 compute/memory，其他切分未测；该扩展是事后结构延伸，未预注册，但继承预注册的平坦对照。

\section{6. 范围、切割与推广空间}\label{scope-differentiation-and-extension-space}

\textbf{这是评分基础，不是解法。} 本文交付的是一个"节点信息增量如何计算、如何验伪"的协议。它不解决 reward hacking、不恢复真实目标、不验证任何目标扩展机制。

\textbf{与既有工作的切割。} 信息增益、roofline、PID 是既有工具。Specification Self-Correction（Gallego 2025, arXiv:2507.18742）在推理时做自然语言的规范批判与改写，不使用信息增益、树结构、条件熵或量化评分，也不训练权重。本文与它的关系是\emph{前置}：本文提供"给节点打分"的可测量基础，不涉及规范修订这一步。

\begin{longtable}[]{@{}
  >{\raggedright\arraybackslash}p{(\linewidth - 12\tabcolsep) * \real{0.1429}}
  >{\raggedright\arraybackslash}p{(\linewidth - 12\tabcolsep) * \real{0.1429}}
  >{\raggedright\arraybackslash}p{(\linewidth - 12\tabcolsep) * \real{0.1429}}
  >{\raggedright\arraybackslash}p{(\linewidth - 12\tabcolsep) * \real{0.1429}}
  >{\raggedright\arraybackslash}p{(\linewidth - 12\tabcolsep) * \real{0.1429}}
  >{\raggedright\arraybackslash}p{(\linewidth - 12\tabcolsep) * \real{0.1429}}
  >{\raggedright\arraybackslash}p{(\linewidth - 12\tabcolsep) * \real{0.1429}}@{}}
\toprule\noalign{}
\begin{minipage}[b]{\linewidth}\raggedright
工作
\end{minipage} & \begin{minipage}[b]{\linewidth}\raggedright
输入
\end{minipage} & \begin{minipage}[b]{\linewidth}\raggedright
输出
\end{minipage} & \begin{minipage}[b]{\linewidth}\raggedright
闭式可算
\end{minipage} & \begin{minipage}[b]{\linewidth}\raggedright
精确零负控制
\end{minipage} & \begin{minipage}[b]{\linewidth}\raggedright
预注册
\end{minipage} & \begin{minipage}[b]{\linewidth}\raggedright
涉及规范修订
\end{minipage} \\
\midrule\noalign{}
\endhead
\bottomrule\noalign{}
\endlastfoot
ID3 信息增益 & 数据分布 & 分裂属性 & 否（需数据） & 无 & 无 & 否 \\
Roofline \citep{roofline2009} & 算力/带宽 spec & 瓶颈区间 & 是 & 无 & 无 & 否 \\
PID (Williams \& Beer) & 多源分布 & 独有/冗余/协同 & 否（需分布） & 无 & 无 & 否 \\
SSC (2507.18742) & rubric + 输出 & 改写后的规范 & 否（推理时生成） & 无 & 无 & \textbf{是} \\
本文 & 硬件 spec + 先验 & 节点增量 $H(D|S)$ & \textbf{是} & \textbf{是（=0）} & \textbf{是} & 否（前置） \\
\end{longtable}

这张表的作用是划界而非争先：本文不复制任何一行的功能。本文提供的是上表未覆盖的一格——节点信息增量可闭式计算，且可被精确零负控制推翻。

\textbf{推广空间（未实现）。} 本文的树是最简单的示例结构。更复杂的版本可以引入新数据的贝叶斯更新、多假设约束、以及最小描述长度准则下的结构选择。这些方向在本稿中不实现、不主张结果；此处仅声明结构的可扩展性边界。

\section{7. 局限}\label{limitations}

\begin{itemize}
\tightlist
\item
  翻转率与占比移动是\textbf{可解释分数}，不是互信息本身；严格量是 $H(D|S)$。
\item
  4090 与 A100 无原生 FP4，其 FP4 速率以 $0.5\times$ FP16 代理，故 4090 的极端比率部分是代理产物。
\item
  工作负载先验为均匀 Dirichlet，未建模真实低精度为主的 LLM 推理；平坦对照零化与结构观察在构造上先验鲁棒，而绝对翻转率会随偏置先验变化。
\item
  单调性不成立（A100 的 shift 大于 5090）；四个平台不足以支撑 scaling law。
\item
  树扩展是事后的、4 点的、固定切分的；层内类型对照是设计层论证，不是测得的臂。
\item
  单一域：四个平台都是 GPU，锚点为手算，未使用实测 trace。外部公开数据的 spot-check 是未来工作。
\end{itemize}

\section{8. 结论}\label{conclusion}

我们示范了一个给分解树节点打分的闭式、可负控制、可预注册的协议。在 GPU 瓶颈诊断这个玩具域上，浅层判断 43--56\% 的情形判错；深层增量为 0.17--0.23 的时间占比移动和 0.71--0.77 bits 的深层独有信息，且在平坦对照下精确归零。结论限定为该评分基础在此域的可计算性与可验伪性：它可计算、可复现、可被自身的负控制推翻。它既不宣称普适性，也不宣称解决任何对齐问题。

\begin{center}\rule{0.5\linewidth}{0.5pt}\end{center}


\bibliography{../submission-staging/references}
\bibliographystyle{tmlr}

\end{document}
